
\section*{5. Relations entre \(\operatorname{Gr}^\bullet(X)\) et \(H^{2*}(X,\underline{Z}_\ell(*))\)}

\textbf{5.1.} L'exemple donné dans 4.8 n'exclut pas qu'on puisse peut-être trouver, pour un schéma $X$ quelconque, un homomorphisme d'anneaux
\[
 g^\bullet:\operatorname{Gr}^\bullet(X)\longrightarrow H^{2*}(X,\underline{Z}_\ell(*)),
\]
tels que pour tout Module localement libre $\underline{E}$ sur $X$, $g^i$ transforme la classe de Chern $c_i(\underline{E})$ de Exp. V en la classe de Chern cohomologique de SGA 5 VII. Bien entendu, ici $\ell$ désigne un nombre premier distinct des caractéristiques résiduelles de $X$. Comme l'anneau $\operatorname{Gr}^\bullet(X)$ est engendré par les $c^i(\underline{E})$, on voit en tous cas que $g^\bullet$ est uniquement déterminé par la condition d'être compatible avec les classes de Chern.

\textbf{5.2.} Utilisant l'homomorphisme de Chern
\[
 K^\bullet(X)\longrightarrow\Ch\bigl(H^{2*}(X,\Zell(*))\bigr),
\]
on en conclut, par passage aux gradués associés respectivement à la filtration du premier membre et la filtration par le degré du second, un homomorphisme de groupes gradués (mais pas d'anneaux !)
\[
 \psi:\operatorname{Gr}^\bullet(X)\longrightarrow H^{2*}(X,\Zell(*)),
\]
qui satisfait à la relation
\[
 \psi(c_i(\underline{E}))=(-1)^{i-1}(i-1)!\,c_i^\ell(\underline{E}),
\]
où $c_i^\ell(\underline{E})$ désigne la classe de Chern $\ell$-adique. On en conclut, par division à partir des $\psi^i$, un homomorphisme d'anneaux :
\[
 g_\Q:\operatorname{Gr}^\bullet(X)\longrightarrow
 H^{2*}(X,\Qell(*))=
 H^{2*}(X,\Zell(*))\otimes_{Z_\ell}\Qell,
\]
transformant classe de Chern en classes de Chern; et tel que l'image du premier membre soit contenue dans celle de $H^2(X,\Zell(*))$; par suite, la question de l'existence de $g^\bullet$ se résoud affirmativement pourvu qu'on remplace $H^{2*}(X,\Zell(*))$ par son quotient par son sous-groupe de torsion. Le premier cas douteux pour définir les composantes $g^i$ de l'hypothètique homomorphisme $g$ est celui de $g^3$, on n'arrive à priori qu'à définir $2g^3=\psi^3$. La réponse à la question d'existence de $g^3$ serait négative si par exemple on pouvait trouver trois faisceaux inversibles $\underline{L}_1,\underline{L}_2,\underline{L}_3$ sur $X$ satisfaisant la relation
\[
 (\underline{L}_1-1)(\underline{L}_2-1)(\underline{L}_3-1)=0\quad\text{dans }K^\bullet(X),
\]
